\chapter{Roteiro de entrevista com gestores}
\label{ap:roteiro}

Este apêndice apresenta o roteiro semiestruturado a ser aplicado em entrevistas com gestores (proprietários, sócios ou gerentes seniores) dos estabelecimentos selecionados como casos. A entrevista tem duração estimada de 45 a 60 minutos, é gravada em áudio mediante assinatura do Termo de Consentimento Livre e Esclarecido (Apêndice~\ref{ap:termos}) e segue uma estrutura modular composta por blocos comuns a todos os perfis e blocos específicos aplicáveis conforme o perfil do estabelecimento: A (pré-2020 que adotou cardápio digital), B (pré-2020 que resistiu à adoção) ou C (pós-2020 nascido digital). As perguntas-âncora abaixo orientam a condução; o entrevistador deve aprofundar respostas com perguntas auxiliares do tipo ``por quê?'' e ``poderia dar um exemplo concreto?''.

\noindent\textbf{Bloco 0, Caracterização} \emph{(todos os perfis)}
\begin{enumerate}
	\item Conte brevemente a história do estabelecimento: quando abriu, qual é o conceito e qual é o perfil de cliente.
	\item Qual é o seu papel no estabelecimento e há quanto tempo está nesse papel?
	\item Quantas mesas o estabelecimento possui, quantos funcionários trabalham regularmente e qual é a faixa de tíquete médio?
	\item Hoje, como é o cardápio em uso (físico, digital ou híbrido)? Há quanto tempo nesse formato?
\end{enumerate}

\noindent\textbf{Bloco 1, Trajetória de adoção} \emph{(perfis A e B)}
\begin{enumerate}\setcounter{enumi}{4}
	\item Como era a operação de cardápio antes de 2020?
	\item O que mudou (ou não) durante e após a pandemia em relação ao cardápio?
	\item \emph{(Se adotou)} O que motivou a decisão de digitalizar? Quem tomou essa decisão? Foi planejada ou emergencial?
	\item \emph{(Se adotou)} Como foi o processo de transição? Quanto tempo levou e quais foram as principais dificuldades?
	\item \emph{(Se resistiu)} Por que optaram por não digitalizar? Já consideraram a adoção? O que mudaria essa decisão?
\end{enumerate}

\noindent\textbf{Bloco 2, Operação atual com cardápio digital} \emph{(perfis A e C)}
\begin{enumerate}\setcounter{enumi}{9}
	\item Qual modalidade de autoatendimento é utilizada: acesso pelo smartphone do cliente (via QR Code), tablet disponibilizado na mesa, totem em quiosque ou alguma combinação? O que motivou essa escolha?
	\item Qual ferramenta ou fornecedor de cardápio digital é utilizado? Por que essa escolha? Já houve troca?
	\item Como funciona a atualização do cardápio (preços, novos itens, esgotamento)? Quem é responsável?
	\item Qual é o custo aproximado (mensalidade, taxas, equipamentos) e qual a percepção sobre o custo-benefício?
	\item Que problemas técnicos são enfrentados com maior frequência (conexão, falha em escanear o QR Code ou abrir o cardápio, tela táctil de tablet ou totem não responsiva, \emph{layout} confuso, dificuldade de uso por clientes)?
	\item Como a digitalização afetou o tempo de atendimento, o tíquete médio, o número de garçons e o retrabalho na cozinha?
\end{enumerate}

\noindent\textbf{Bloco 3, Resistência e barreiras} \emph{(perfil B)}
\begin{enumerate}\setcounter{enumi}{14}
	\item Que receios você tem em relação ao cardápio digital (custos, dependência de fornecedor, perda da experiência presencial)?
	\item Como você imagina que seus clientes reagiriam? Já houve algum tipo de teste com cardápio digital?
	\item O que precisaria acontecer para vocês adotarem cardápio digital?
\end{enumerate}

\noindent\textbf{Bloco 4, Percepção do cliente} \emph{(todos os perfis)}
\begin{enumerate}\setcounter{enumi}{17}
	\item Que tipo de feedback (positivo e negativo) vocês recebem dos clientes sobre o cardápio?
	\item Vocês observam diferenças de comportamento entre grupos distintos (idade, turistas \emph{versus} locais, grupos grandes \emph{versus} casais)?
	\item Há clientes que pedem o cardápio físico mesmo havendo a versão digital? Com que frequência isso ocorre?
	\item Vocês têm acesso a dados de uso do cardápio digital (visualizações, itens mais consultados)? Esses dados são utilizados para alguma decisão?
\end{enumerate}

\noindent\textbf{Bloco 5, Olhar para o futuro} \emph{(todos os perfis)}
\begin{enumerate}\setcounter{enumi}{21}
	\item Em sua avaliação, onde essa digitalização chegará no seu estabelecimento daqui a cinco anos?
	\item Se pudesse mudar uma única coisa na solução de cardápio digital que existe hoje, o que seria?
	\item Há algo importante sobre esse tema que não foi abordado nesta entrevista?
\end{enumerate}

\noindent\emph{Observação metodológica:} antes da coleta oficial será realizado um piloto da entrevista com pelo menos um respondente do setor (não computado entre os casos), com o objetivo de revisar a clareza das perguntas, identificar redundâncias e ajustar a duração estimada.
